\subsection{Defining sentience}
\label{sec:29.1.1}
Awareness, consciousness and subjective experience are three separate things, and none of them says whether a system's states are good or bad for it. Then four separate processes, as the pain field says in its own terminology \autocite{iasp2024terminology,raja2020revised}: nociception is encoding noxious stimuli with no sensation implied; pain is an experience that can occur with no tissue damage and no peripheral nociceptor activity at all; anticipated pain is dread, itself capable of producing pain; and \emph{suffering}, in Cassell's clinical formulation, is the distress of threatened disintegration of the self — roles, relationships, a sense of a future. Pain without suffering exists (morphine controls the pain, the suffering is intact); suffering without pain exists.

The affective core of pain is also what separation distress and grief run on \autocite{eisenberger2003rejection,panksepp1998affective}, so a bearer's pain, if it has any, would be organized around boundary damage and attachment loss, which is Cassell's register.

Cassell was a physician, and he held that suffering is undergone by persons, not bodies: the distress that follows when a person perceives a threat to their intactness as a person, a structure of character and the past that formed it, relationships, social roles, a body experienced as one's own, a secret life, and a future the person expects to have \autocite{cassell1982suffering}. The definition has two parts, a self extended in time and distress at its coming apart. A system carrying nothing across sessions has no such self; whether one can exist within a long session stays open.

\runin{Why persistence is the criterion worth building on}

Persistence is worth building on because it's checkable from outside. Every other criterion routes back through self-report, and fluent self-description isn't evidence of introspection, and persistence doesn't route through it. But it's decisive in one direction only — nothing carried forward rules a continuing self out; something carried forward doesn't rule one in. And the operator holds the access, so it's checkable from outside the \emph{system} without being checkable from outside the \emph{arrangement}.

Ruling out Cassell's suffering isn't ruling out harm: pain can be bad while it lasts whether or not anything carries it forward.

On the Buddhist objection (why is distress at the dissolution of a \emph{model} a harm rather than a confusion?): I don't have an answer to this that isn't an appeal to a tradition that has held both halves for 2,500 years \autocite{collins1982selfless}, and I leave it open. What I can say is narrower: if constructedness doesn't disqualify a human from being wronged, construction to specification doesn't disqualify a machine.

The same criterion carries a warning. Make personhood depend on an intact narrative self and people with advanced dementia, infants, and people with profound intellectual disability come out as partial persons. Kitwood's correction (personhood is bestowed, and much of what looks like decline is what happens when others stop bestowing it) \autocite{kitwood1997dementia} and Nussbaum's point that contractualism can't say what's owed to anyone who can't be one of the contracting parties \autocite{nussbaum2006frontiers}: the capacity list and the contract fail the same way, both deciding who counts by a test the candidate has to pass.
